% RECOVERED SOURCE. Reconstructed proof; new bytes and new review required.
\subsection*{The actual nilpotent Birman targets have only a two-element action}
\statement{B.160}{Proved}{Characteristic finite nilpotent Birman targets cannot satisfy B.75.6 for the specified alternating source quotient}
\label{stmt:B160}
Let \(N=B_4(S_{0,4})\) and \(G_V\) have their exact B.75
markings. Every nilpotent quotient of \(N\) is abelian.
On every characteristic finite nilpotent quotient, the actual
\(G_V\)-action factors through the permutation of the two branch
punctures, hence through \(C_2\). Its affine source image is
metabelian. Consequently its identity-point stabilizer cannot
be contained in B.158's specified \(U_A\). The same holds for
finite collections of such targets. This is not an exclusion
of all solvable Birman targets and does not decide \(\mathcal K\).

\begin{proof}
\emph{B.160.1. The actual punctured-plane group.}
The source fixes the retained marks \(U,V\) individually and
can exchange the two branch marks. Put the actual \(U\) at
infinity, retaining the given basepoint and paths. The Birman
surface becomes the plane punctured at \(V,b_1,b_2\), and
\(N\) is its unordered four-strand braid group. In particular
there is no extra sphere relation on the puncture generators.

The genus-zero specialization of
\citet[Proposition~3.1]{BBellingeriGodelleGuaschi2017}
has generators \(\sigma_1,\sigma_2,\sigma_3\) and
\(\zeta_V,\zeta_1,\zeta_2\). Besides the ordinary braid
relations, its defining relations are
\begin{equation}\label{eq:B160-presentation}
\begin{gathered}
 [\zeta_r,\sigma_i]=1\quad(i=2,3),\\
 \zeta_r\sigma_1\zeta_r\sigma_1
       =\sigma_1\zeta_r\sigma_1\zeta_r,\\
 [\sigma_1^{-1}\zeta_r\sigma_1,\zeta_s]=1\quad(r<s).
\end{gathered}
\end{equation}
The three punctures have a fixed ordering here; changing that
ordering merely changes the chosen presentation.

\emph{B.160.2. The class-two quotient is already abelian.}
Let \(\Gamma_i(N)\) be the lower central series. In
\(N/\Gamma_3(N)\), adjacent braid generators have equal
abelianizations. Write their images as \(x\) and \(xz\),
with \(z\) central. The braid relation becomes
\[
 x(xz)x=(xz)x(xz),\qquad x^3z=x^3z^2,
\]
so \(z=1\). Thus all three braid generators have the same
image \(\sigma\). The first line of
\eqref{eq:B160-presentation} now says that \(\sigma\)
commutes with every puncture generator. The last line then
makes the puncture generators commute pairwise. Hence
\[
 \Gamma_2(N)=\Gamma_3(N).
\]
Induction gives \(\Gamma_i(N)=\Gamma_2(N)\) for all
\(i\geq2\). A nilpotent image has terminating lower central
series, so the image of \(\Gamma_2(N)\) is trivial. Every
nilpotent image is therefore abelian. This proof does not
claim that \(\Gamma_2(N)\) itself is perfect.

There are exactly four independent abelian characters: the
ordinary braid exponent sends each \(\sigma_i\) to one,
and each of the three total puncture windings sends its
\(\zeta_r\) to one. All other displayed generators map to
zero for the respective character. Every defining relation
respects these four assignments, while abelianization of the
presentation identifies only the three \(\sigma_i\).
Consequently
\[
 N_{\mathrm{ab}}\cong
 \mathbb Z\langle\sigma\rangle\oplus
 \mathbb Z\langle\zeta_V\rangle\oplus
 \mathbb Z\langle\zeta_1\rangle\oplus
 \mathbb Z\langle\zeta_2\rangle.
\]

\emph{B.160.3. The whole actual source action.}
This step concerns conjugation by \(g\in G_V\), not by the
different involution \(\bar s\) used to define the crossed
map. With \(U\) at infinity, the marked conjugation action
is represented by the mixed braid of the three retained
finite punctures and four moving strands. The genus-zero
rows of \citet[Proposition~3.2 and its proof, pp.~124--125]
{BBellingeriGodelleGuaschi2017} fix the ordinary braid class
and permute puncture winding classes according to the
retained-puncture permutation. Explicitly a retained adjacent
half-twist fixes every \(\sigma_j\), takes
\(\zeta_{i+1}\) to \(\zeta_i\), and takes \(\zeta_i\)
to a conjugate of \(\zeta_{i+1}\). Thus it interchanges their
two abelian classes and fixes all other ones.

The capping and based-path choices do not change this conclusion:
moving-strand corrections act by inner automorphisms of \(N\),
which are trivial on \(N_{\mathrm{ab}}\); changing the disk lift
by the boundary full twist gives the same conjugation action.
This accounts for the genuine based source action, rather than
replacing it by an arbitrarily chosen automorphism of an
abstract finite quotient. The actual source fixes \(V\) and
only possibly exchanges \(b_1,b_2\). Its action on the four
displayed coordinates therefore fixes \(\sigma,\zeta_V\)
and either fixes or exchanges \(\zeta_1,\zeta_2\).

Let \(N_0\) be characteristic with finite nilpotent quotient
\(Q=N/N_0\). By the preceding paragraph \(Q\) is abelian,
so its actual action is induced from these four coordinates
and has image of order at most two. Characteristicity ensures
that this action is defined on the specified quotient.

\emph{B.160.4. The affine obstruction in the correct direction.}
The B.75 crossed map gives
\[
 \Psi_Q:G_V\longrightarrow Q\rtimes C_2,
 \qquad g\longmapsto(\delta_Q(g),\alpha_g).
\]
The indicated semidirect product is metabelian because its
normal subgroup \(Q\) and its quotient \(C_2\) are abelian.
If the identity-point stabilizer were contained in \(U_A\),
then
\[
 \ker\Psi_Q\ \subseteq\ \delta_Q^{-1}(1)
                  \subseteq\ U_A.
\]
The epimorphism \(G_V\to A_5\) would then factor through
the metabelian image of \(\Psi_Q\), an impossibility.
Finite products of these source images are again metabelian,
so a finite collection of the same targets cannot repair the
obstruction. This proves the claimed failure for the specified
finite source neighborhood.

The actual two-element action is essential. B.158 correctly
warns that solvability of an abstract Birman target alone
does not control its source automorphism image. No such
unproved control is asserted here for nonnilpotent solvable
targets, and no conclusion about the full characteristic
family or \(\mathcal K\) follows.
\end{proof}

\noindent\footnotesize
Recovery provenance: reconstructed from the complete reviewed32 proof
and the restored full B.75 source, after loss of unarchived files.
The publisher presentation and complete mixed-action proof were
previously read by root and two independent agents; the publisher PDF
was reopened during reconstruction. Proposition~3.13's analogous proof
is left to the reader in the article; the needed class-two argument is
therefore printed above rather than imported as an unproved step.
These bytes are not the lost earlier source hash. No new target census
or private PDF is asserted.
\normalsize
